%=====================================================================
\chapter{Confidential Computing and Emerging Hardware}\label{ch:confidential}
%=====================================================================
\locator{6}{Part VI --- Emerging Hardware and Practice}

\begin{outcomes}
By the end of this chapter you will be able to:
\begin{tight}
  \item \textbf{State} how the trust model inverts, and \textbf{say} what the
        guest stops trusting.
  \item \textbf{Explain} memory encryption with per-guest keys, and
        \textbf{say} what it does and does not protect against.
  \item \textbf{Describe} remote attestation and \textbf{explain} why
        encryption without it is nearly worthless.
  \item \textbf{Compare} AMD SEV-SNP, Intel TDX and ARM CCA by maturity and
        availability as of late 2026.
  \item \textbf{Decide} whether a stated workload justifies a confidential VM,
        and \textbf{state} the cost.
\end{tight}
\end{outcomes}

\begin{prereq}
\chref{security} (the attack surface, and in particular who is trusted) and
\chref{hardware} (VMX, the VMCS and extended page tables, which are what this
hardware modifies). \chref{cloud}'s shared responsibility model is the
commercial arrangement this chapter challenges.
\end{prereq}

\begin{keyterms}
confidential computing $\bullet$ trusted execution environment $\bullet$ TEE
$\bullet$ trust model $\bullet$ memory encryption $\bullet$ SEV $\bullet$
SEV-ES $\bullet$ SEV-SNP $\bullet$ TDX $\bullet$ trust domain $\bullet$ ARM CCA
$\bullet$ realm $\bullet$ SGX $\bullet$ enclave $\bullet$ attestation $\bullet$
attestation report $\bullet$ measurement $\bullet$ chain of trust $\bullet$
confidential container $\bullet$ CXL $\bullet$ composable infrastructure
\end{keyterms}

\section{The Trust Model Inverts}

Every chapter so far has assumed the hypervisor is trustworthy. It schedules
guests, allocates their memory and mediates their devices, and a guest has no
choice but to rely on it. \chref{security} worried about an attacker
\emph{reaching} the hypervisor; it never questioned the hypervisor itself.

Confidential computing questions it. The guest is given a way to run on
hardware it does not own, operated by people it has not met, \emph{without
trusting them with its data}.

This matters commercially rather than philosophically. A hospital will not put
patient records in a public cloud if the provider's administrators can read
them. A bank will not run a model on hardware whose operator could copy the
weights. The cloud's answer has always been contracts and audits; confidential
computing offers a technical answer instead.

\begin{definitionbox}[Trusted Execution Environment]
A \term{TEE} is an execution context whose memory and registers are protected
from everything outside it --- including the hypervisor, the host operating
system, the firmware and anyone with physical access.

\smallskip
For a confidential VM this means:
\begin{tight}
  \item memory is \textbf{encrypted} with a key the guest's own hardware context
        holds and the host cannot read;
  \item register state is protected across every VM exit, so the hypervisor
        cannot inspect the guest's execution;
  \item the guest can \textbf{prove} what it is --- which hardware, which
        firmware, which image --- to a remote party.
\end{tight}

\smallskip
The hypervisor still \emph{schedules} the guest, allocates it memory and gives
it devices. It simply cannot see inside it.
\end{definitionbox}

\dsfig{%
\begin{tikzpicture}[font=\scriptsize,
  bx/.style={rounded corners=2pt, line width=0.9pt, align=center,
             inner sep=3pt, minimum height=0.6cm, font=\scriptsize},
  t/.style={bx, draw=greenhl, fill=greenhl!14, text=greenhl!45!black},
  u/.style={bx, draw=accent, fill=accent!10, text=accent!70!black},
  hd/.style={font=\scriptsize\bfseries, align=center},
  note/.style={font=\scriptsize\itshape, align=center}
]
  % --- conventional ---
  \node[hd, text=primary] at (3.0,4.3) {Conventional};
  \node[t, minimum width=2.6cm] at (1.6,3.5) {guest};
  \node[t, minimum width=2.6cm] at (4.4,3.5) {guest};
  \node[t, minimum width=5.6cm] at (3.0,2.7) {hypervisor};
  \node[t, minimum width=5.6cm] at (3.0,1.9) {firmware};
  \node[t, minimum width=5.6cm] at (3.0,1.1) {hardware};
  \node[t, minimum width=5.6cm] at (3.0,0.3) {the operator};
  \node[note, text=greenhl!45!black] at (3.0,-0.4)
       {the guest trusts \emph{all} of this};

  % --- confidential ---
  \node[hd, text=primary] at (10.6,4.3) {Confidential};
  \node[t, minimum width=2.6cm] at (9.2,3.5) {guest};
  \node[u, minimum width=2.6cm, draw=neutral!55, fill=neutral!10,
        text=neutral!45!black] at (12.0,3.5) {another guest};
  \node[u, minimum width=5.6cm] at (10.6,2.7) {hypervisor};
  \node[u, minimum width=5.6cm] at (10.6,1.9) {firmware};
  \node[t, minimum width=5.6cm] at (10.6,1.1) {hardware \ \
        {\scriptsize\mdseries CPU + its root of trust}};
  \node[u, minimum width=5.6cm] at (10.6,0.3) {the operator};
  \node[note, text=accent!70!black] at (10.6,-0.4)
       {red is \emph{untrusted} --- and still runs the machine};

  \draw[decorate, decoration={brace, amplitude=4pt}, greenhl, line width=0.9pt]
        (13.6,3.8) -- (13.6,0.8);
  \node[note, text=greenhl!45!black, anchor=west, align=left] at (13.9,2.3)
       {the trusted base\\
        shrinks to the\\
        guest and the\\
        silicon};
\end{tikzpicture}}%
{The inversion. The hypervisor still schedules the guest and allocates its
memory; it simply loses the ability to read it. Everything red is in the
attack path and outside the trust boundary.}{trust-inversion}

\section{What the Hardware Does}

\begin{conceptbox}[Memory Encryption, and Why Encryption Alone Is Not Enough]
The memory controller encrypts and decrypts transparently with a key tied to
the guest, held in the processor and never exposed to software. Host memory
dumped from outside yields ciphertext.

\smallskip
That defeats the straightforward attacks --- a hypervisor reading guest pages,
a memory image taken from a running host, cold-boot attacks on the DIMMs. AMD's
first \term{SEV} (2017) did this much, and researchers quickly showed it was
not enough.

\smallskip
The reason is that a hostile hypervisor does more than read. It can
\emph{remap} pages --- give the guest a different physical page than it had ---
\emph{replay} old encrypted pages, or \emph{alias} two guest addresses to one
physical page. The data stays encrypted throughout and the guest is still
compromised.

\smallskip
\term{SEV-ES} (2018) added register-state protection across VM exits.
\term{SEV-SNP} (2020) added the missing piece: \emph{integrity}. A reverse map
table records which guest owns each physical page, and the hardware refuses a
mapping that contradicts it. Intel's \term{TDX} was designed with equivalent
protections from the start.

\smallskip
The lesson generalises: confidentiality without integrity is a weaker property
than it sounds, and the first generation of almost every such system learns
this the same way.
\end{conceptbox}

\begin{definitionbox}[Attestation]
\term{Attestation} is the guest proving what it is to a remote party before that
party entrusts it with anything.

\smallskip
The hardware measures the guest's initial state --- firmware, kernel, initial
memory image --- and produces an \term{attestation report} containing those
measurements, the processor's security version, and a signature chaining to a
key the silicon vendor certified.

\smallskip
The verifier checks the signature, confirms the hardware is genuine and
up to date, and compares the measurement against what it expected. Only then
does it release the decryption key, the model weights or the database
credentials.

\smallskip
\textbf{Encryption without attestation protects nothing useful.} If you cannot
verify what is running, a hostile operator simply boots a modified image that
asks you politely for the data --- and that image's memory is beautifully
encrypted while it exfiltrates. Attestation is the part that makes the
encryption meaningful, and it is the part deployments most often skip.
\end{definitionbox}

\section{Where the Platforms Stand}

Dated October 2026. This is the fastest-moving material in the book; check
before relying on any row.

\rowcolors{2}{white}{lightbg}
\begin{longtable}{@{}L{2.6cm}L{4.0cm}L{4.0cm}L{4.1cm}@{}}
\toprule
\rowcolor{primary!15}
& \textbf{AMD SEV-SNP} & \textbf{Intel TDX} & \textbf{ARM CCA} \\
\midrule
\endhead
Unit protected & The whole VM & The whole VM (a \term{trust domain}) & The whole
VM (a \term{realm}) \\
Introduced & SEV 2017, SNP 2020 & 2023 & Announced 2021 \\
Maturity & The most deployed; generally available across major clouds &
Production confidential VMs on Azure and Google Cloud, with newer machine
families still arriving in preview during 2026 & Specification mature,
\emph{silicon and cloud availability still limited}; previews rather than
general availability \\
Watch for & Attestation report format changed to v4 with a firmware update ---
parsers written for v3 break & A TDX firmware vulnerability was disclosed and
patched in early 2026; keep firmware current & Expect availability to change
fastest of the three \\
\bottomrule
\end{longtable}

\begin{alertbox}[The predecessor that is not this: Intel SGX]
SGX protected an \emph{enclave} --- a region inside one process --- rather than
a whole virtual machine. The application had to be rewritten and split, with
the sensitive part inside the enclave and everything else outside, across a
narrow and awkward boundary. Enclave memory was also small, in the low hundreds
of megabytes on most parts.

\smallskip
It saw real use and it did not become general infrastructure, for the obvious
reason: nobody wants to re-architect an application to run it somewhere.

\smallskip
TDX, SEV-SNP and CCA all learned from that. They protect an \emph{unmodified
virtual machine}, so an existing workload moves in without being rewritten ---
which is the same bargain \chref{vmm} struck when it insisted on running
unmodified guests, made again two decades later.
\end{alertbox}

\begin{worked}[Does This Workload Justify a Confidential VM?]
A hospital wants to run analytics over patient records in a public cloud. The
compliance team objects that the provider's staff could read the data. Decide
whether confidential computing answers the objection, and what it costs.

\smallskip
\textbf{Step 1 --- what is actually being claimed?} Not that the provider
\emph{would}, but that they \emph{could}, and that the hospital cannot
demonstrate otherwise to its regulator. The requirement is therefore
evidentiary, not merely technical.

\smallskip
\textbf{Step 2 --- what a confidential VM gives.}
\begin{tight}
  \item memory encrypted with a key the host cannot read;
  \item integrity protection against remapping and replay (SNP or TDX);
  \item an attestation report the hospital can verify before releasing the
        database key.
\end{tight}
The provider's administrators can still stop the machine, and still see how
much memory and processor it uses. They cannot read the records.

\smallskip
\textbf{Step 3 --- what it does \emph{not} give.} Four gaps worth stating to
the compliance team before they believe too much:
\begin{tight}
  \item \textbf{Availability.} The operator can still destroy or detain the
        machine. Confidentiality and integrity; not availability.
  \item \textbf{Side channels.} Timing and cache attacks
        (\chref{security}) are not fully addressed, and some published attacks
        target exactly these systems.
  \item \textbf{Your own bugs.} A flaw in the application inside the enclave
        leaks the data the normal way.
  \item \textbf{Data in use elsewhere.} The moment records leave for an
        unprotected service, the guarantee ends. The boundary must cover the
        whole path.
\end{tight}

\smallskip
\textbf{Step 4 --- the cost.} Performance overhead is workload-dependent and
generally modest --- commonly in the region of 2--8\% for ordinary server
workloads, more where I/O is heavy, because encrypted guests use bounce buffers
for device transfers. Instance choice is restricted to supported families, and
some features --- live migration (\chref{migration}), certain memory management
--- may be unavailable or restricted.

\smallskip
The larger cost is operational: attestation must be built, verified and
maintained. Somebody must decide what measurement is expected, store it, check
it, and update it with every image change. \textbf{A confidential VM whose
attestation nobody verifies is a normal VM with a performance penalty.}

\smallskip
\textbf{Step 5 --- the recommendation.} Yes, with a condition. Use a
confidential VM \emph{and} build the attestation flow, so the database key is
released only to a verified measurement. Record the expected measurement under
change control. If the project cannot commit to maintaining that, the honest
advice is to keep the data on premises, because the alternative is a control
that exists only on a slide.
\end{worked}

\section{Composable Infrastructure and CXL}

One more direction the hardware is moving, included because the 2017 outline
asked about the future of virtualization and this is a plausible part of it.

Every host in this book is a fixed box: so many cores, so much memory, bought
together and discarded together. A host with spare memory cannot lend it to the
host next to it, so every machine carries headroom nobody else can use ---
stranded capacity, at scale.

\term{CXL} --- Compute Express Link --- allows memory to be attached to a fabric
rather than to a socket, so a pool of memory can be allocated to whichever host
needs it. The ambition, \term{composable infrastructure}, is to assemble a
machine from pooled components for a workload and return them afterwards.

\begin{conceptbox}[Why This Is the Same Idea as Chapter 1]
Pooled memory, allocated on demand and returned, is \emph{virtualization of the
memory bus} --- multiplexing, by the definition of \chref{what}. The motive is
the one from the opening pages: capacity bought and left idle because it could
not be shared.

\smallskip
The obstacle is the one the whole book keeps meeting: latency. Fabric-attached
memory is slower than socket-attached memory, as NUMA's remote node is slower
than its local one (\chref{cpumem}) --- only more so. Early deployments are
therefore tiered, with hot pages local and cold pages pooled, which is exactly
the memory-reclamation ladder of \chref{cpumem} wearing different clothes.

\smallskip
Treat this section as a direction, not a product. The useful thing to carry
away is that you already have the framework to evaluate it when a vendor
arrives: ask what is multiplexed, what the lie costs, and who maintains it.
\end{conceptbox}

\begin{pitfall}
\begin{tight}
  \item \textbf{Deploying encryption without attestation.} You have hidden the
        data from an attacker who was not going to read it, and not from one
        who boots a modified image.
  \item \textbf{Believing the first generation is enough.} Plain SEV lacked
        integrity and was broken by remapping. Use SEV-SNP or TDX.
  \item \textbf{Assuming availability is protected.} It is not. The operator
        can still stop your machine.
  \item \textbf{Assuming side channels are solved.} They are not, and published
        attacks target these systems specifically.
  \item \textbf{Hard-coding an attestation report format.} AMD's move to v4
        broke parsers written for v3.
  \item \textbf{Forgetting the rest of the path.} The guarantee ends where the
        data leaves the TEE.
  \item \textbf{Planning on ARM CCA availability.} As of late 2026 it is
        previews, not general availability. Check before designing around it.
\end{tight}
\end{pitfall}

\begin{lab}[Attestation, Verified]
Inspects a confidential VM and verifies an attestation report. Where no
suitable processor is available, the captured reports in \texttt{code/} make
the verification exercise identical. Expect thirty-five minutes.

\begin{lstlisting}[language=bash]
#!/usr/bin/env bash
# ch21_attest.sh -- is this guest confidential, and can it prove it?
set -euo pipefail

# --- 1. Does this processor support any of the three? -----------------
grep -o -m1 -E '\<(sev|sev_es|sev_snp|tdx_guest)\>' /proc/cpuinfo || \
  echo "no TEE features on this host"
ls /sys/firmware/ 2>/dev/null | grep -iE 'tdx|sev' || true
dmesg | grep -iE 'SEV-SNP|TDX|Memory Encryption' | head -5 || true

# --- 2. From inside a guest: am I confidential? -----------------------
# Each platform exposes itself differently; check all three.
[ -e /dev/sev-guest ]  && echo "AMD SEV-SNP guest device present"
[ -e /dev/tdx_guest ]  && echo "Intel TDX guest device present"
systemd-detect-virt
# On a normal guest, none of these exist. That is the point of step 4.

# --- 3. Ask the hardware for an attestation report. -------------------
# snpguest (AMD) or trustee/tdx tools (Intel). A 64-byte nonce from the
# verifier goes in; a signed report comes out.
NONCE=$(openssl rand -hex 32)
if [ -e /dev/sev-guest ] && command -v snpguest >/dev/null; then
  snpguest report report.bin request.txt --random
  snpguest display report report.bin | head -30
else
  echo "no TEE here -- using the captured report in code/"
  cp "$(dirname "$0")/ch21_report.txt" report.txt 2>/dev/null || true
  sed -n '1,30p' report.txt 2>/dev/null || true
fi

# --- 4. The three fields that matter, and why. ------------------------
cat <<'NOTES'
  measurement    the hash of the guest's initial state. Compare this
                 with what you built; if it differs, something else
                 is running and you must NOT release the key.
  report_data    your nonce, echoed back. Proves the report is fresh
                 and not a recording of an earlier genuine boot.
  signature +    chains to a key the silicon vendor certified. Proves
  cert chain     the report came from real hardware, not from software
                 pretending very convincingly.
NOTES

# --- 5. Verify the chain. This is the step people skip. ---------------
if command -v snpguest >/dev/null; then
  snpguest fetch ca-der milan ./certs           # vendor certificates
  snpguest fetch vcek-der milan ./certs report.bin
  snpguest verify certs ./certs
  snpguest verify attestation ./certs report.bin && \
    echo "VERIFIED: genuine hardware, measurement as reported"
else
  echo "verification exercise: compare the measurement in report.txt"
  echo "with code/ch21_expected_measurement.txt and decide."
fi

# --- 6. Now the attack this defends against. --------------------------
# A hostile operator boots a DIFFERENT image that asks for your key.
# Its memory is perfectly encrypted while it exfiltrates. The only
# thing that stops it is step 5 refusing the measurement.
diff <(grep -i measurement report.txt 2>/dev/null || echo a) \
     <(cat "$(dirname "$0")/ch21_expected_measurement.txt" 2>/dev/null \
       || echo b) >/dev/null \
  && echo "measurement matches -- release the key" \
  || echo "measurement DIFFERS -- refuse. This is the whole mechanism."

# --- 7. What it costs: the same benchmark, confidential or not. ------
# On a cloud that offers both, run this on a confidential and a normal
# instance of the same family and compare.
sysbench cpu --threads=4 --time=20 run 2>/dev/null | grep -E 'events/s' \
  || echo "install sysbench to measure the overhead"
dd if=/dev/zero of=/tmp/t bs=1M count=2048 oflag=direct 2>&1 | tail -1
rm -f /tmp/t
# I/O usually shows the larger penalty: encrypted guests bounce device
# transfers through shared buffers.
\end{lstlisting}

\textbf{What to notice.} Step 6 is the chapter. Encryption was working perfectly
in both the honest and the hostile case; the only thing that distinguished them
was a hash comparison that somebody had to perform. Skip step 5 and you have
bought a performance penalty.
\end{lab}

\begin{checkpoint}
\begin{tightnum}
  \item Explain why memory encryption without integrity protection was
        insufficient, naming one concrete attack it failed to stop.
  \item A team deploys confidential VMs but releases the database key from a
        configuration file at boot, without checking any report. State exactly
        what they are protected against and what they are not.
  \item Name two properties confidential computing does \emph{not} provide, and
        say why each matters to a compliance argument.
\end{tightnum}
\end{checkpoint}

\begin{chaptersummary}
\begin{tight}
  \item Confidential computing inverts the trust model: the guest stops
        trusting the hypervisor, the firmware and the operator, and trusts only
        itself and the silicon.
  \item A TEE encrypts guest memory with a key the host cannot read, protects
        register state across exits, and lets the guest prove what it is.
  \item Encryption alone was not enough: a hostile hypervisor can remap, replay
        and alias pages. SEV-SNP and TDX add integrity.
  \item Attestation is what makes the encryption meaningful. A confidential VM
        whose attestation nobody verifies is a normal VM with a performance
        penalty.
  \item As of late 2026 SEV-SNP is the most widely deployed, TDX is in
        production on major clouds with newer families still arriving, and ARM
        CCA remains limited in availability.
  \item It does not protect availability, does not solve side channels, and
        does not help with your own bugs. Overhead is commonly 2--8\%, more for
        I/O.
  \item CXL and composable infrastructure are the same multiplexing idea
        applied to the memory bus, with latency as the familiar obstacle.
\end{tight}
\end{chaptersummary}

\begin{reviewq}
  \item \textbf{(MCQ)} SEV-SNP added, over SEV-ES: (a)~larger enclaves;
        (b)~integrity protection against page remapping and replay;
        (c)~live migration; (d)~support for containers.
  \item \textbf{(MCQ)} Without attestation, memory encryption fails to stop:
        (a)~a cold-boot attack; (b)~a hypervisor reading guest pages; (c)~an
        operator booting a modified image that requests the data; (d)~a
        physical memory dump.
  \item \textbf{(Short)} Explain the three fields of an attestation report and
        what each one proves.
  \item \textbf{(Short)} Why did Intel SGX not become general infrastructure,
        and what did the VM-scoped designs do differently?
  \item \textbf{(Short)} State three things confidential computing does not
        protect.
  \item \textbf{(Applied)} A bank wants to run a proprietary model on a public
        cloud without the provider being able to copy the weights. Set out the
        design, the verification step that makes it real, what it costs, and
        the residual risks you would record.
  \item \textbf{(Applied)} Using the framework of \chref{what} --- what is
        multiplexed, what does the lie cost, who maintains it --- evaluate CXL
        pooled memory as a vendor would present it to you, and write the three
        questions you would ask.
\end{reviewq}
